\chapter{Sorting and Conversion}
\label{ch:sorting}

\section*{The Sorting Mode\index{modes!Sorting Mode|textbf}}

\begin{versebox}
The teaching, its basis, and its ground ---\\
sort what's shared from what's unique.\\
Two kinds of sutta, two kinds of path:\\
know which one you seek.
\end{versebox}

\noindent Not all suttas are doing the same job. This seems obvious once you say it, but the implications are enormous. The Sorting Mode is the tool that tells you which kind of sutta you are reading --- and therefore how to read it.\footnote{Pali: \textit{``Tattha katamo vibhattihāro? `Dhammañca padaṭṭhānaṃ bhūmiñcā'ti.''} The mnemonic fragment is ``the teaching, its basis, and its ground'' (\textit{dhammaṃ ca padaṭṭhānaṃ bhūmiñ ca}). The word \textit{vibhatti} means ``sorting, classification, analysis'' --- the mode that divides the undivided into categories.}

There are two kinds of sutta. Suttas for cultivation, and suttas for penetration.\footnote{Pali: \textit{``Dve suttāni vāsanābhāgiyañca nibbedhabhāgiyañca.''} \textit{Vāsanābhāgiya} suttas are ``cultivation-sharing'' --- they build good habits, prepare the ground. \textit{Nibbedhabhāgiya} suttas are ``penetration-sharing'' --- they cut through to direct understanding.}

A cultivation sutta\index{sutta types!cultivation} talks about generosity, about virtue, about the heavenly realms, about the drawbacks of sensual pleasure and the benefits of renunciation. It builds the foundation. A penetration sutta\index{sutta types!penetration} illuminates the Four Noble Truths. It breaks through.

Here is the critical difference: in a cultivation sutta, there is no direct understanding, no path, no fruit. In a penetration sutta, there \textit{is} direct understanding, there \textit{is} the path, there \textit{is} the fruit.\footnote{Pali: \textit{``Vāsanābhāgiyaṃ nāma suttaṃ dānakathā sīlakathā saggakathā kāmānaṃ ādīnavo nekkhamme ānisaṃso'ti. \ldots Nibbedhabhāgiyaṃ nāma suttaṃ yā catusaccappakāsanā, vāsanābhāgiye sutte natthi pajānanā, natthi maggo, natthi phalaṃ. Nibbedhabhāgiye sutte atthi pajānanā, atthi maggo, atthi phalaṃ.''} Five topics of cultivation suttas: generosity (\textit{dāna}), virtue (\textit{sīla}), heaven (\textit{sagga}), the danger of sense-pleasures (\textit{kāmānaṃ ādīnava}), and the benefit of renunciation (\textit{nekkhamme ānisaṃsa}). This is the standard ``graduated teaching'' (\textit{anupubbikathā}) that the Buddha would give before turning to the Four Noble Truths.}

Both matter. You cannot penetrate without first cultivating. But confusing the two --- reading a cultivation sutta as if it were a penetration sutta, or vice versa --- leads to misunderstanding.

Two kinds of sutta lead to two kinds of practice: the merit-oriented practice and the fruit-oriented practice. And two kinds of virtue: restraint-virtue (not doing what harms) and abandonment-virtue (letting go of what binds). The person standing on restraint-virtue, practicing through a cultivation sutta, lives the holy life at one level. The person standing on abandonment-virtue, practicing through a penetration sutta, lives the holy life at another.\footnote{Pali: \textit{``Dve paṭipadā puññabhāgiyā ca phalabhāgiyā ca. Dve sīlāni saṃvarasīlañca pahānasīlañca.''} The two-by-two structure: cultivation suttas / merit-oriented practice / restraint-virtue vs.\ penetration suttas / fruit-oriented practice / abandonment-virtue. These are not separate paths but stages of the same journey.}

\section*{What's Shared and What's Not}

The Sorting Mode then tackles a subtle question: which defilements do different types of practitioners have in common, and which are unique to their level?\footnote{Pali: \textit{``Tattha katame dhammā sādhāraṇā?''} ``What qualities are shared?'' The Netti distinguishes two types of commonality: \textit{nāmasādhāraṇā} (shared in name) and \textit{vatthusādhāraṇā} (shared in basis). Two practitioners might share the \textit{name} ``trainee'' (\textit{sekha}) while being at very different levels.}

An ordinary person and a stream-enterer\index{stream-entry} share sensual desire and ill-will in common. Both can experience these --- the stream-enterer has not yet eradicated them, only weakened the view-based fetters. An ordinary person and a non-returner share the higher fetters in common --- subtle craving for existence, restlessness, conceit, ignorance. And whatever worldly attainment a noble disciple enters, that attainment is shared in common with those who have not yet been freed from lust.\footnote{Pali: \textit{``Puthujjanassa sotāpannassa ca kāmarāgabyāpādā sādhāraṇā, puthujjanassa anāgāmissa ca uddhaṃbhāgiyā saṃyojanā sādhāraṇā, yaṃ kiñci ariyasāvako lokiyaṃ samāpattiṃ samāpajjati, sabbā sā avītarāgehi sādhāraṇā.''} A revealing passage: even noble disciples, when they enter worldly meditative attainments, are in a state ``shared'' with those who haven't been freed from lust. The attainment itself is the same; the underlying freedom is different.}

But here is what is \textit{not} shared: the underlying nature of their attainment. The eighth person --- the one on the very cusp of stream-entry --- shares the \textit{name} ``stream-enterer'' with the stream-enterer, but their underlying nature is completely different. All trainees share the \textit{name} ``trainee,'' but the nature of a stream-enterer's training and a once-returner's training are worlds apart.\footnote{Pali: \textit{``Aṭṭhamakassa sotāpannassa ca kāmarāgabyāpādā sādhāraṇā dhammatā asādhāraṇā, aṭṭhamakassa anāgāmissa ca uddhambhāgiyā saṃyojanā sādhāraṇā dhammatā asādhāraṇā. Sabbesaṃ sekkhānaṃ nāmaṃ sādhāraṇaṃ dhammatā asādhāraṇā.''} Name shared, nature not shared. This distinction matters enormously for interpretation: when a sutta says ``the trainee,'' the Sorting Mode helps you determine \textit{which level} of trainee is meant. The ``eighth person'' (\textit{aṭṭhamaka}) is the person on the supramundane path to stream-entry --- not yet a stream-enterer but no longer an ordinary person.}

The Sorting Mode provides a comprehensive map of these relationships --- which practices serve as the basis for which stages, how generosity relates to hearing the teaching, how virtue relates to wise reflection, how development relates to right view, how living in a suitable place supports both seclusion and concentration, and how the Teacher's excellence serves to inspire those without faith and deepen those who already have it.\footnote{Pali: \textit{``Dānamayaṃ puññakiriyavatthu parato ghosassa sādhāraṇaṃ padaṭṭhānaṃ \ldots sappurisūpanissayo tiṇṇañca aveccappasādānaṃ samathassa ca sādhāraṇaṃ padaṭṭhānaṃ \ldots satthusampadā appasannānañca pasādāya pasannānañca bhiyyobhāvāya sādhāraṇaṃ padaṭṭhānaṃ.''} The Netti maps an elaborate web of relationships between merit-bases, practice-supports, and stages of the path. The key insight is that everything is connected: generosity supports hearing the teaching; virtue supports wise reflection; development supports penetrative wisdom; suitable environment supports both concentration and seclusion; good companions support verified confidence and calm; right self-direction supports conscience and insight; abandoning the unwholesome supports investigation and concentration; the Teacher's excellence supports faith for the faithless and deepens faith for the faithful.}

The point is not to memorize the map but to grasp the principle: when you read a sutta, you need to know which level it is addressing. A teaching aimed at stream-enterers will not make sense if you read it as addressed to ordinary people, and vice versa. The Sorting Mode keeps you from reading the wrong manual for your stage.\footnote{Pali: \textit{``Tenāha āyasmā mahākaccāyano `dhammañca padaṭṭhānan'ti.''}}

\ornament

\section*{The Conversion Mode\index{modes!Conversion Mode|textbf}}

\begin{versebox}
Wholesome and unwholesome states ---\\
when one arises, one must fall.\\
Every gain is also loss;\\
each choice converts them all.
\end{versebox}

\noindent The Conversion Mode is about opposites. When you pick up one thing, you automatically put down another.\footnote{Pali: \textit{``Tattha katamo parivattano hāro? `Kusalākusale dhamme'ti.''} The mnemonic fragment is ``wholesome and unwholesome states'' (\textit{kusalākusale dhamme}). The word \textit{parivattana} means ``turning around, converting'' --- when one quality arises, its opposite is automatically converted into absence.}

Think of a light switch. You cannot have the light on \textit{and} off simultaneously. The Conversion Mode says: every right factor, when it arises, \textit{automatically} converts its wrong counterpart into absence. And every wrong factor that arises converts its right counterpart.

When right view arises in a person, wrong view is overcome. And whatever unwholesome states would have arisen because of wrong view --- those too are overcome. And because of right view, many wholesome states arise and reach fulfillment. The same for right intention, right speech, right action, right livelihood, right effort, right mindfulness, right concentration, right liberation, and right knowledge-and-vision of liberation.\footnote{Pali: \textit{``Sammādiṭṭhissa purisapuggalassa micchādiṭṭhi nijjiṇṇā bhavati. Ye cassa micchādiṭṭhipaccayā uppajjeyyuṃ aneke pāpakā akusalā dhammā, te cassa nijjiṇṇā honti. Sammādiṭṭhipaccayā cassa aneke kusalā dhammā sambhavanti \ldots Evaṃ sammāvācassa \ldots sammāvimuttiñāṇadassanassa \ldots te cassa bhāvanāpāripūriṃ gacchanti.''} The Netti applies this conversion pattern to all ten factors of the ``extended'' path (the standard eight plus right liberation and right knowledge-and-vision). Each right factor has a triple effect: (1) it overcomes its wrong counterpart, (2) it overcomes all the unwholesome states that would have arisen from that wrong counterpart, and (3) it generates wholesome states that reach fulfillment. \textbf{A note on our translation}: the original Pali states this pattern ten times, once for each factor. We have summarized the pattern because the logic is identical in every case --- the full repetition is preserved in this note.}

The same principle works for conduct. When you refrain from killing, killing is abandoned. When you refrain from stealing, stealing is abandoned. When you live in celibacy, unchastity is abandoned. When you speak truthfully, lying is abandoned. And so on through each of the ten courses of action.\footnote{Pali: \textit{``Yassa vā pāṇātipātā paṭiviratassa pāṇātipāto pahīno hoti. Adinnādānā paṭiviratassa adinnādānaṃ pahīnaṃ hoti. Brahmacārissa abrahmacariyaṃ pahīnaṃ hoti \ldots Sammādiṭṭhissa micchādiṭṭhi pahīnā hoti.''} Ten conversions: refraining from (1) killing, (2) stealing, (3) unchastity, (4) lying, (5) divisive speech, (6) harsh speech, (7) idle chatter → each abandons its opposite. Plus (8) non-covetousness abandons covetousness, (9) non-ill-will abandons ill-will, (10) right view abandons wrong view.}

\section*{The Logical Consequence}

The Netti then draws out a sharp logical point. Whoever criticizes the Noble Eightfold Path --- they are, whether they realize it or not, praising right view's opposite. By attacking right view, they logically commit themselves to honoring wrong view. By attacking right intention, they honor wrong intention. And so on down the line.\footnote{Pali: \textit{``Ye ca kho keci ariyaṃ aṭṭhaṅgikaṃ maggaṃ garahanti, nesaṃ sandiṭṭhikā sahadhammikā gārayhā vādānuvādā āgacchanti. Sammādiṭṭhiñca te bhavanto dhammaṃ garahanti. Tena hi ye micchādiṭṭhikā, tesaṃ bhavantānaṃ pujjā ca pāsaṃsā ca.''} A devastating argument: criticism of the path is logically equivalent to praise of its opposite. You cannot reject right view without implicitly endorsing wrong view. You cannot reject right speech without implicitly endorsing wrong speech. The Conversion Mode makes this logical implication explicit.}

And whoever says ``Sense-pleasures should be enjoyed! Indulged in! Cultivated! Pursued!'' --- for such people, restraint from sense-pleasures is wrong. Whoever says ``Self-mortification is the teaching'' --- for such people, the path that actually leads out is wrong. Whoever says ``The teaching is painful'' --- for such people, the teaching that brings happiness is wrong.\footnote{Pali: \textit{``Ye ca kho keci evamāhaṃsu `bhuñjitabbā kāmā, paribhuñjitabbā kāmā \ldots' Kāmehi veramaṇī tesaṃ adhammo. Ye vā pana keci evamāhaṃsu `attakilamathānuyogo dhammo'ti. Niyyāniko tesaṃ dhammo adhammo. Ye ca kho keci evamāhaṃsu `dukkho dhammo'ti. Sukho tesaṃ dhammo adhammo.''} Three opposing views and their conversions: (1) hedonists → restraint is their non-teaching; (2) ascetics → the path that leads out is their non-teaching; (3) pessimists → the happiness of the path is their non-teaching. Each wrong position automatically rejects its correct counterpart.}

The deepest application is in meditation itself. When a monk contemplates the unattractive, the perception of beauty is abandoned. When he contemplates the painful, the perception of pleasure is abandoned. When he contemplates the impermanent, the perception of permanence is abandoned. When he contemplates not-self, the perception of self is abandoned.\footnote{Pali: \textit{``Yathā vā pana bhikkhuno sabbasaṅkhāresu asubhānupassino viharato subhasaññā pahīyanti. Dukkhānupassino viharato sukhasaññā pahīyanti. Aniccānupassino viharato niccasaññā pahīyanti. Anattānupassino viharato attasaññā pahīyanti.''} The four contemplations and their four conversions: unattractive → beauty-perception abandoned; painful → pleasure-perception abandoned; impermanent → permanence-perception abandoned; not-self → self-perception abandoned. These are the four ``antidotes'' to the four distortions (\textit{vipallāsa}) discussed in Chapter~\ref{ch:basis}.}

Whatever you embrace, its opposite is automatically rejected. Whatever teaching you accept and practice, the opposite teaching is automatically overcome. This is the Conversion Mode: every choice is a double choice, every gain a double gain.\footnote{Pali: \textit{``Yaṃ yaṃ vā pana dhammaṃ rocayati vā upagacchati vā, tassa tassa dhammassa yo paṭipakkho, svassa aniṭṭhato ajjhāpanno bhavati. Tenāha āyasmā mahākaccāyano `kusalākusaladhamme'ti.''} The closing formula. The principle is universal: \textit{``Whatever quality one favors or approaches, the opposite of that quality is automatically, through what is unwanted, brought upon oneself.''} Every step toward light is a step away from darkness. Every step away from darkness is a step toward light.}
